\begin{helper}[Terminal recoloring gives a shorter bad path]
Let \((\varphi,T)\) be a minimum bad pair with
\(T=(y_0,e_1,y_1,\ldots,e_n,y_n)\), \(n\geq2\), and \(\rho(T)=1\).
Let \(p<n\) satisfy \(p+1=n\), and suppose the color \(\beta\) is
missing at both \(y_p=y_{n-1}\) and \(y_n\).  Then \(p\geq1\), and
there are a color \(\delta\), an index \(j<p\), and a coloring
\(\varphi^*\) with the following properties:
\begin{enumerate}
\item the old value of \(e_n\) is \(\delta\), and \(\delta\) is missing
at \(y_j\) under \(\varphi\);
\item \(\varphi^*\) is obtained by changing exactly \(e_n\) from
\(\delta\) to \(\beta\); and
\item the exact prefix \(T_{n-1}\) is a bad pair under \(\varphi^*\)
and has smaller Tashkinov rank than \(T\).
\end{enumerate}
\end{helper}
\begin{proof}
The equations \(p+1=n\) and \(n\geq2\) give \(p\geq1\).  Since
\(\rho(T)=1\), the defining equations in
\noderef{TashTreeSuffixPath} say that the last exact edge, which is the
edge in position \(p\) in zero-based Lean indexing, has endpoint pair
\(\{y_p,y_n\}\).

This last edge is distinct from \(e_1=e_0\).  Indeed, its position \(p\)
is positive, while the edge-position map of \noderef{TreeSequence} is
injective.  The \noderef{KTriple} inside
\noderef{TashTreeBadPair}, obtained from the first conjunct of
\noderef{TashTreeMinimumBadPair}, supplies
\noderef{IsOneEdgeUncolored}.  Its second clause therefore supplies a
color \(\delta\) with \(\varphi(e_n)=\delta\).

Because the last edge has positive displayed index, the
\noderef{TashkinovTree} clause of badness supplies an earlier displayed
vertex \(y_j\), \(j\leq p\), at which \(\delta\) is missing.  In fact
\(j<p\).  If \(j=p\), then \(e_n\) itself would be an incident edge of
color \(\delta\) at \(y_p\), putting \(\delta\) in
\noderef{PresentColors} there and excluding it from
\noderef{MissingColors}.  This contradicts the witness assertion.
Thus \(j<p\), establishing the old color and the strictly earlier
witness required in the statement.

Define the promised coloring by the exact-edge formula
\[
 \varphi^*(e)=
 \begin{cases}
   \beta,&e=e_n,\\
   \varphi(e),&e\ne e_n.
 \end{cases}
\]
It is proper.  A pair of distinct adjacent edges not containing \(e_n\)
is unchanged and is handled by \noderef{IsProperEdgeColoring} for
\(\varphi\).  If the pair contains \(e_n\), its other edge shares one of
the endpoints \(y_p,y_n\), by \noderef{AdjacentEdges} and the endpoint
equation above.  Since \(\beta\) is missing at both endpoints, the
definitions \noderef{PresentColors} and \noderef{MissingColors} show
that the other edge does not have value \(\beta\).  Thus recoloring
\(e_n\) with \(\beta\) creates no conflict.  The uniquely uncolored edge
is still \(e_0\): \(e_n\ne e_0\) and one colored value was replaced by
another.  Therefore \noderef{IsOneEdgeUncolored} is preserved.  All
other clauses of \noderef{KTriple} concern only the graph, \(k\), and
the fixed chromatic and degree data, so \((G,e_0,\varphi^*)\) is again a
\(k\)-triple.

Let \(P\) be the exact \noderef{TreeSequencePrefix} of length \(p\).
It is an initial segment in the sense of
\noderef{TreeSequenceInitialSegment}, and \(p<n\).  Every exact edge of
\(P\) differs from \(e_n\) by edge-position injectivity, so its value is
unchanged by the definition of \(\varphi^*\).  Fix such a displayed edge
from the second position onward.  Its old witness from
\noderef{TashkinovTree} occurs strictly before that edge, hence before
position \(p\).  The only endpoints whose incident color sets can change
are \(y_p,y_n\), so this old witness is neither endpoint and retains the
same missing color.  Thus each old witness remains valid and \(P\) is a
Tashkinov tree under \(\varphi^*\).

The vertices \(y_j,y_p\) show that \(P\) is not elementary under
\(\varphi^*\).  At \(y_j\), the missingness of \(\delta\) is unchanged,
because \(j<p\) makes \(y_j\) distinct from both endpoints of \(e_n\).
At \(y_p\), properness of the old coloring implies that \(e_n\) was the
only incident edge with value \(\delta\): a second one would be a
distinct adjacent edge with the same color, contrary to
\noderef{IsProperEdgeColoring}.  Replacing the value of \(e_n\) therefore
removes the last incident \(\delta\)-edge, so \(\delta\) is missing at
\(y_p\) under \(\varphi^*\).  The two vertices are distinct because
\(j<p\) and the vertex map of \noderef{TreeSequence} is injective.  By
\noderef{ColoredVertexSetElementary} and
\noderef{TreeSequenceVertexSet}, the prefix is non-elementary.

The new \(k\)-triple, the verified Tashkinov condition, and this
non-elementariness are precisely \noderef{TashTreeBadPair} for
\((\varphi^*,P)\).  Finally \(P\) is a proper exact initial segment of
\(T\), so \noderef{TashTreeProperInitialSegmentRankLess} gives
\noderef{TashTreeRankLess} of \(P\) below \(T\).  This proves every
asserted part of the terminal replacement.
\end{proof}
